\documentclass[11pt,a4paper]{article}

\usepackage[utf8]{inputenc}
\usepackage[T1]{fontenc}
\usepackage[english]{babel}
\usepackage{lmodern}
\usepackage{microtype}

\usepackage{amsmath,amssymb}
\usepackage{braket}

\usepackage[margin=25mm,headheight=15pt,headsep=8mm]{geometry}
\usepackage{fancyhdr}
\usepackage{needspace}
\usepackage{titling}
\setlength{\droptitle}{-12mm}
\setlength{\parindent}{0pt}
\setlength{\parskip}{0.5\baselineskip plus 1pt minus 1pt}
\setlength{\jot}{5pt}
\setlength{\emergencystretch}{2em}
\raggedbottom
\clubpenalty=10000
\widowpenalty=10000
\displaywidowpenalty=10000
\pagestyle{fancy}
\fancyhf{}
\fancyhead[L]{\small\nouppercase{Superconducting qubit readout}}
\fancyhead[R]{\small Kacper Kłos}
\fancyfoot[C]{\thepage}
\renewcommand{\headrulewidth}{0.3pt}

\usepackage[numbers,sort&compress]{natbib}
\usepackage[hidelinks]{hyperref}
\hypersetup{
  pdftitle={Superconducting qubit readout: an open quantum system perspective},
  pdfauthor={Kacper Klos}
}
\numberwithin{equation}{section}

\title{Superconducting qubit readout:\\[0.3em]
  \large An open quantum system perspective}
\author{Kacper Kłos}
\date{}

\begin{document}
\maketitle
\thispagestyle{plain}

\section{Transmon--resonator system}

We begin our analysis of the transmon--resonator system with the Hamiltonian
given in Eq.~(31) of Blais et al.~\cite{Blais2021}:
\begin{equation}
  \hat H = 4E_C(\hat n+\hat n_r)^2-E_J\cos\varphi +\sum_m\hbar\omega_m\hat a_m^\dagger\hat a_m,
\end{equation}
where $\hat n_r=\sum_m\hat n_m$.

We separate this Hamiltonian into three contributions: the transmon
Hamiltonian $\hat H_t$, the resonator Hamiltonian $\hat H_r$, and their
interaction $\hat H_i$:
\begin{align}
  \hat H   &= \hat H_t+\hat H_r+\hat H_i, \\
  \hat H_t &= 4E_C\hat n^2-E_J\cos\varphi, \\
  \hat H_r &= 4E_C\hat n_r^2 +\sum_m\hbar\omega_m\hat a_m^\dagger\hat a_m, \\
  \hat H_i &= 8E_C\hat n\hat n_r.
\end{align}
We examine each contribution in turn.

\subsection{Transmon Hamiltonian}

To determine how $\cos\hat\varphi$ acts on a charge state $\ket{n}$,
we first relate the operators $\hat\varphi$ and $\hat n$.
Using $\hat n=\hat Q/(2e)$, $\hat\varphi=(2e/\hbar)\hat\Phi$, and
$[\hat\Phi,\hat Q]=i$, we obtain
\begin{equation}
  [\hat\varphi,\hat n] =\frac{1}{\hbar}[\hat\Phi,\hat Q]=i.
\end{equation}
We can then determine the action of $e^{i\hat\varphi}$ on $\ket{n}$.
Hadamard's lemma gives
\begin{equation}
  \begin{aligned}
    e^{i\hat\varphi}\hat n e^{-i\hat\varphi}
      &=\hat n+[i\hat\varphi,\hat n]
        +\frac12[i\hat\varphi,[i\hat\varphi,\hat n]]\\
      &=\hat n-1.
  \end{aligned}
\end{equation}
It follows that
\begin{equation}
  \hat n\bigl(e^{-i\hat\varphi}\ket{n}\bigr) =(n-1)e^{-i\hat\varphi}\ket{n}.
\end{equation}
Thus, $e^{-i\hat\varphi}\ket{n}$ is an eigenstate of $\hat n$ with
eigenvalue $n-1$. Applying the same reasoning to
$e^{i\hat\varphi}\ket{n}$ gives
$e^{\pm i\hat\varphi}\ket{n}=\ket{n\pm1}$.
Together with
$\cos\varphi=\tfrac12(e^{i\hat\varphi}+e^{-i\hat\varphi})$,
this determines the action of the transmon Hamiltonian:
\begin{equation}
  \hat H_t\ket{n} =4E_Cn^2\ket{n} -\frac{E_J}{2}\bigl(\ket{n+1}+\ket{n-1}\bigr).
\end{equation}

\subsection{Resonator Hamiltonian}

The relations needed to expand the resonator Hamiltonian are
\begin{equation}
  \begin{aligned}
    \hat n_m &=\frac{C_g}{2e\,C_m}\hat Q_m,\\
    \hat Q_m &=i\sqrt{\frac{\hbar C_m\omega_m}{2}}
      (\hat a_m^\dagger-\hat a_m).
  \end{aligned}
\end{equation}
Consequently,
\begin{equation}
  \hat n_m=i\eta_m(\hat a_m^\dagger-\hat a_m),
  \qquad
  \eta_m=\frac{C_g}{2e}\sqrt{\frac{\hbar\omega_m}{2C_m}}.
\end{equation}
We now adopt a single-mode approximation, retaining only one resonator
mode. Substituting these expressions into $\hat H_r$ and omitting
constant terms gives
\begin{equation}
  \hat H_r =(8E_C\eta^2+\hbar\omega_0)\hat a^\dagger\hat a -4E_C\eta^2\bigl((\hat a^\dagger)^2+\hat a^2\bigr).
\end{equation}
To simplify this Hamiltonian, we introduce a Bogoliubov transformation:
\begin{align}
  \hat{\tilde a}
    &=\cosh(r)\hat a-\sinh(r)\hat a^\dagger, \\
  \hat a
    &=\cosh(r)\hat{\tilde a}+\sinh(r)\hat{\tilde a}^\dagger, \\
  [\hat{\tilde a},\hat{\tilde a}^\dagger]
    &=\cosh^2(r)-\sinh^2(r)=1.
\end{align}
Substituting into the operator products appearing in $\hat H_r$ yields
\begin{equation}
  \begin{aligned}
    \hat a^\dagger\hat a
      &=(\cosh(r)\hat{\tilde a}^\dagger+\sinh(r)\hat{\tilde a}) (\cosh(r)\hat{\tilde a}+\sinh(r)\hat{\tilde a}^\dagger) \\
      &=\frac12\sinh(2r) \bigl((\hat{\tilde a}^\dagger)^2+\hat{\tilde a}^2\bigr)+\cosh(2r)\hat{\tilde a}^\dagger\hat{\tilde a}
          +\sinh^2(r), \\
    \hat a^2
      &=\cosh^2(r)\hat{\tilde a}^2+\sinh^2(r)(\hat{\tilde a}^\dagger)^2+\sinh(2r)\left(\hat{\tilde a}^\dagger\hat{\tilde a}+\frac12\right), \\
    (\hat{\tilde a}^\dagger)^2+\hat{\tilde a}^2
      &=\cosh(2r)\bigl((\hat a^\dagger)^2+\hat a^2\bigr)+\sinh(2r) \bigl(2\hat{\tilde a}^\dagger\hat{\tilde a}+1\bigr).
  \end{aligned}
\end{equation}
For compactness, define
\begin{equation}
  A=8E_C\eta^2+\hbar\omega_0,
  \qquad B=4E_C\eta^2.
\end{equation}
Substituting the preceding expressions into $\hat H_r$ gives
\begin{equation}
  \begin{aligned}
    \hat H_r
      &=\bigl(A\cosh(2r)-2B\sinh(2r)\bigr)
        \hat{\tilde a}^\dagger\hat{\tilde a}+\left(\frac A2\sinh(2r)-B\cosh(2r)\right) \bigl((\hat{\tilde a}^\dagger)^2+\hat{\tilde a}^2\bigr)+\text{const.}
  \end{aligned}
\end{equation}
We choose $r$ so that the terms quadratic in the creation or annihilation
operator cancel:
\begin{equation}
  \begin{aligned}
    \tanh(2r)
      &=\frac{2B}{A}=\frac{8E_C\eta^2}{\hbar\omega_r},\\[3pt]
    r 
      &=\frac14\ln\left(\frac{A+2B}{A-2B}\right)
      =\frac14\ln\left(1+\frac{16E_C\eta^2}{\hbar\omega_0}\right).
  \end{aligned}
\end{equation}
In the final Hamiltonian, we denote the transformed operator
$\hat{\tilde a}$ by $\hat a$ for simplicity. The transformation must
nevertheless be accounted for when expressing the interaction term.
The resonator Hamiltonian becomes
\begin{equation}
  \hat H_r=\hbar\omega\hat a^\dagger\hat a+\text{const.},
\end{equation}
where
\begin{equation}
  \begin{aligned}
    \hbar\omega
      &=A\cosh(2r)-2B\sinh(2r)\\
      &=\sqrt{A^2-4B^2}
       =\omega_0\sqrt{1+\frac{16E_C\eta^2}{\hbar\omega_0}}.
  \end{aligned}
\end{equation}
Here we have used standard hyperbolic-function identities to obtain
$\omega$. In this form, the action of $\hat H_r$ on its eigenstates
$\ket{j}$ is immediate.

\subsection{Interaction Hamiltonian}

In addition to the approximation introduced for the resonator, we must
account for the change in the creation and annihilation operators under
the Bogoliubov transformation:
\begin{equation}
  \begin{aligned}
    \hat a^\dagger-\hat a
      &=\bigl(\cosh(r)-\sinh(r)\bigr) (\hat{\tilde a}^\dagger-\hat{\tilde a})\\
      &=e^{-r}(\hat{\tilde a}^\dagger-\hat{\tilde a}).
  \end{aligned}
\end{equation}
As in the preceding subsection, we use $\hat a$ to denote the
transformed annihilation operator in the final expression:
\begin{equation}
  \hat H_i
  =8iE_C\eta\hat n e^{-r}(\hat a^\dagger-\hat a)
  =ig(\hat a^\dagger-\hat a).
\end{equation}

\subsection{Total Hamiltonian}

Combining the contributions above gives
\begin{equation}
    \hat H
      =4E_C\hat n^2 -\frac{E_J}{2}(e^{i\varphi}+e^{-i\varphi})+\hbar\omega\hat a^\dagger\hat a+ig(\hat a^\dagger-\hat a).
\end{equation}
We express this Hamiltonian in the basis of Cooper-pair charge states
associated with the Josephson junction and dressed-photon states of the
resonator.

\newpage

\bibliographystyle{unsrtnat}
\bibliography{references}

\end{document}
